% Figure: pharmacokinetic concentration-time curve (Bateman function)
% and the cumulative area under the curve (AUC). The concentration rises
% as the drug absorbs and falls as it clears; the shaded area to a given
% time is the exposure integral, and the total area is the AUC that sets
% the dose-exposure relationship.
\begin{figure}[htbp]
\centering
\begin{tikzpicture}
\begin{axis}[
  width=11cm, height=6.2cm,
  xlabel={time $t$ (hours)}, ylabel={plasma concentration $C(t)$ (mg/L)},
  xmin=0, xmax=12, ymin=0, ymax=6.5,
  domain=0:12, samples=200,
  axis lines=left, grid=both,
  major grid style={engblack!12}, font=\small,
  legend pos=north east, legend style={font=\footnotesize},
]
% Bateman: C(t) = 10*(exp(-0.4 t) - exp(-1.5 t)), ka=1.5, ke=0.4
% shade AUC up to t = 5
\addplot[draw=none, fill=engblue!15, domain=0:5]
  {10*(exp(-0.4*x) - exp(-1.5*x))} \closedcycle;
\addplot[engblue, very thick]
  {10*(exp(-0.4*x) - exp(-1.5*x))};
\addlegendentry{$C(t)$ (Bateman)}
% peak marker: tmax = ln(ka/ke)/(ka-ke) = ln(3.75)/1.1 = 1.202
\addplot[engred, only marks, mark=*, mark size=1.6pt]
  coordinates {(1.202,4.10)};
\node[font=\scriptsize, anchor=south west, engred]
  at (axis cs:1.3,4.15) {$C_{\max}$ at $t_{\max}$};
\node[font=\scriptsize, engblue] at (axis cs:2.6,1.1)
  {$\int_{0}^{5} C\,\dd t$};
\end{axis}
\end{tikzpicture}
\caption[Concentration-time curve and area under the curve]{A plasma
concentration-time curve after a single oral dose, modelled by the
Bateman function $C(t) = A\bigl(e^{-k_{e}t} - e^{-k_{a}t}\bigr)$ with an
absorption rate $k_{a}$ and an elimination rate $k_{e}$. The
concentration climbs while absorption outpaces clearance, peaks at
$t_{\max}$, and decays once clearance dominates. The shaded region is
the exposure integral $\int_{0}^{5} C\,\dd t$ accumulated to five hours;
the total area under the curve, $\mathrm{AUC} = \int_{0}^{\infty}
C\,\dd t = A(1/k_{e} - 1/k_{a})$, is the single number a pharmacologist
reads as total drug exposure and against which a dose is titrated. The
peak, the time to peak, and the area are three integrals-and-derivatives
readings of the same curve.}
\label{fig:vol02ch06:drug-auc}
\end{figure}
